\chapter{Đánh giá thực nghiệm}

\section{Thiết lập đo lường}
Tất cả metric RF thuần dùng cùng tập kiểm tra 20\% với \texttt{random\_state=42}. Mô hình là artifact \texttt{rf\_model\_khtn.pkl}, \texttt{rf\_model\_khxh.pkl} có trong repository tại thời điểm đo.

Đánh giá \textbf{hệ lai} đòi hỏi vòng lặp qua từng mẫu và từng ngành để tìm cực đại điểm hybrid.

\section{Kết quả Random Forest thuần}
\begin{table}[htbp]
\centering
\caption{Hiệu năng phân loại trên tập kiểm tra đầy đủ}
\label{tab:mlmetrics}
\begin{tabular}{lccc}
\toprule
Khối & Accuracy & Macro F1 & Số mẫu test \\
\midrule
KHTN & 0{,}9645 & 0{,}9645 & 33\,498 \\
KHXH & 0{,}9797 & 0{,}9797 & 12\,770 \\
\bottomrule
\end{tabular}
\end{table}

Độ chính xác cao phản ánh mối quan hệ mạnh giữa đặc trưng và nhãn; không nên diễn giải như độ chính xác dự báo lựa chọn nghề nghiệp thực tế của con người.

\section{Ma trận nhầm lẫn (RF, tập kiểm tra đầy đủ)}
Bảng \ref{tab:cmkhtn} và \ref{tab:cmkhxh}: hàng là nhãn đúng, cột là nhãn dự đoán (số nguyên mã ngành).

\begin{table}[htbp]
\centering
\caption{Ma trận nhầm lẫn --- KHTN ($n=33\,498$)}
\label{tab:cmkhtn}
\small
\begin{tabular}{lrrrrr}
\toprule
$\downarrow$ thực / dự đoán $\rightarrow$ & 0 & 1 & 2 & 3 & 4 \\
\midrule
0 IT & 6469 & 84 & 3 & 57 & 87 \\
1 Kinh tế & 106 & 6171 & 91 & 93 & 238 \\
2 Y khoa & 2 & 14 & 6527 & 52 & 105 \\
3 Kỹ thuật & 42 & 23 & 47 & 6488 & 100 \\
4 Nông--Lâm--Ngư & 2 & 21 & 6 & 15 & 6655 \\
\bottomrule
\end{tabular}
\end{table}

\begin{table}[htbp]
\centering
\caption{Ma trận nhầm lẫn --- KHXH ($n=12\,770$)}
\label{tab:cmkhxh}
\small
\begin{tabular}{lrrrr}
\toprule
$\downarrow$ thực / dự đoán $\rightarrow$ & 1 & 5 & 6 & 7 \\
\midrule
1 Kinh tế & 3167 & 14 & 1 & 11 \\
5 Sư phạm & 40 & 3132 & 1 & 19 \\
6 Luật & 45 & 23 & 3085 & 40 \\
7 Du lịch & 19 & 25 & 21 & 3127 \\
\bottomrule
\end{tabular}
\end{table}


\section{Kết quả hệ lai trên mẫu con $N_h=2000$}
Với mỗi hàng kiểm tra, nhãn dự đoán hybrid là ngành có điểm $H_i$ lớn nhất sau khi tính cho mọi $i$ trong tập ngành khối.

\begin{table}[htbp]
\centering
\caption{Độ chính xác hybrid ($N_h=2000$ mẫu đầu của tập test)}
\label{tab:hyb}
\begin{tabular}{lcc}
\toprule
Khối & Accuracy & Macro F1 \\
\midrule
KHTN & 0{,}910 & 0{,}911 \\
KHXH & 0{,}977 & 0{,}977 \\
\bottomrule
\end{tabular}
\end{table}

Trên KHTN, hybrid thấp hơn RF khoảng năm điểm phần trăm tuyệt đối: cơ chế VETO và trọng số 40\% KBS làm thay đổi ranh giới quyết định so với cực đại hóa xác suất thuần túy. Trên KHXH, lai ghép gần với RF, cho thấy luật và xác suất đồng thuận nhiều hơn trên miền nhãn này (trong phân phối đã đo).

\section{Bảng so sánh tổng hợp ML và hybrid}
\begin{table}[!htbp]
\centering
\caption{Tóm tắt metric chính (cùng tập test, hybrid trên $N_h=2000$ mẫu đầu)}
\label{tab:comparemlhyb}
\begin{tabular}{lccc}
\toprule
Khối & RF accuracy (full test) & Hybrid accuracy ($N_h{=}2000$) & Ghi chú \\
\midrule
KHTN & 0{,}9645 & 0{,}910 & VETO làm giảm khớp nhãn proxy \\
KHXH & 0{,}9797 & 0{,}977 & Gần tương đương RF \\
\bottomrule
\end{tabular}
\end{table}
\FloatBarrier

\section{Thảo luận định tính và ý nghĩa đồ án}
\begin{itemize}[leftmargin=*]
  \item \textbf{Ý nghĩa độ chính xác cao của RF:} phản ánh mức tương quan giữa heuristic và không gian điểm; không nên báo cáo như ``độ chính xác tư vấn nghề nghiệp thực tế''.
  \item \textbf{Hybrid thấp hơn trên KHTN:} đánh đổi có chủ đích giữa \textit{khớp nhãn tự động} và \textit{ràng buộc tri thức}; phù hợp ứng dụng cần giải thích an toàn.
  \item \textbf{Ma trận nhầm lẫn:} nhầm giữa IT--Kinh tế và Kinh tế--Nông--Lâm gợi ý ranh giới heuristic không tách biệt tuyến tính hoàn toàn; có thể dùng phân tích điểm biên để hiệu đính luật hoặc điều chỉnh trọng số theo ngành.
\end{itemize}

\section{Độ tin cậy thống kê (nhận thức phương pháp)}
Báo cáo không kèm khoảng tin cậy bootstrap cho accuracy vì nhãn là proxy; nếu mở rộng đồ án, có thể bootstrap trên tập kiểm tra với cỡ mẫu cố định để báo cáo khoảng ước lượng.

\section{Kết chương}
RF đạt hiệu năng phân tách rất tốt trên nhãn proxy; hybrid điều chỉnh theo tri thức với chi phí độ chính xác nhãn có thể đo được trên KHTN.
